\chapter*{Resum}
\addcontentsline{toc}{chapter}{Resum}

El present document recull la proposta de disseny i l'anàlisi de requisits de \textit{Dona'm Bauxa}, una aplicació web orientada a la descoberta i difusió de música en directe, artistes i esdeveniments culturals de l'illa de Mallorca. El projecte s'emmarca en l'assignatura de Tecnologia Multimèdia de la Universitat de les Illes Balears i ha estat desenvolupat per Josep Ferriol Font i Dylan Luigi Canning Garcia.

La plataforma proposada neix per respondre a una necessitat real de l'ecosistema cultural mallorquí: la dispersió de la informació sobre concerts i artistes locals en múltiples canals digitals poc integrats entre si. \textit{Dona'm Bauxa} proposa centralitzar aquesta informació en un únic punt d'accés, accessible des de qualsevol dispositiu, amb una interfície moderna, atractiva i arrelada a la identitat visual de l'illa.

Des del punt de vista funcional, el sistema preveu un cercador d'artistes i grups mallorquins amb fitxes detallades de biografia i discografia, una agenda d'esdeveniments actualitzada amb sistema de filtres per gènere, zona geogràfica i dates, i la integració de tecnologies de geolocalització per representar els espais culturals en un mapa interactiu. Quant a les funcionalitats secundàries, s'inclouen una secció de \textit{Favorits}, accés directe a plataformes d'escolta i compra d'entrades, i un espai de notícies sobre l'actualitat musical balear.

Quant als requisits no funcionals, la plataforma estarà dissenyada seguint els principis de \textit{responsive web design}, complirà els estàndards d'accessibilitat WCAG~2.1, garantirà la seguretat de les comunicacions mitjançant HTTPS i serà compatible amb els principals navegadors tant en dispositius mòbils com en escriptori. L'arquitectura preveu la possibilitat d'escalar cap a la resta de les Illes Balears en fases posteriors de desenvolupament.

En definitiva, \textit{Dona'm Bauxa} aspira a convertir-se en el referent digital de la cultura musical mallorquina, posant en valor el talent local i apropant l'oferta d'oci cultural a tothom qui vulgui gaudir de la bauxa de l'illa.

\bigskip

\noindent\textbf{Paraules clau:} aplicació web, música mallorquina, esdeveniments culturals, geolocalització, disseny responsiu, accessibilitat, anàlisi de requisits.
